\section{Related work}
\label{sec:related}

\paragraph{Multimodal crisis assessment and robustness.}
CrisisMMD pairs social-media text and images with humanitarian and damage annotations from seven 2017 disasters \citep{alam2018crisismmd}. Multimodal systems built on this setting combine both channels for humanitarian categorization and post-level damage severity \citep{ofli2020multimodal,agarwal2020crisisdias,shetty2025disaster}. Newer resources broaden the setting to multi-platform hurricane posts annotated for humanitarian impact and information integrity \citep{yao2026mash}, pixel-level earthquake damage in ground-level imagery \citep{huang2026eidseg}, and VLM-agent post-disaster assessment and reporting \citep{chen2024disasteller}. Exact and near duplicates can inflate disaster-image evaluation unless audited explicitly \citep{alam2020benchmarks}, while crisis-domain adversarial training predates instruction-following VLMs \citep{chen2020crisisadv}. We evaluate zero-shot VLMs without model adaptation, focusing on under-triage of initially correct ordinal decisions.

\paragraph{Typographic and multimodal adversarial attacks.}
Surveys distinguish single-modal manipulation, cross-modal attacks, instruction hijacking, and training-time poisoning \citep{jain2026survey}. Optimization-based work trains adversarial images to control VLM outputs \citep{bailey2024imagehijacks} or coordinates image and text perturbations for white- or black-box transfer \citep{zhang2022coattack,yin2023vlattack}. A complementary semantic surface is visible language: rendered text can alter recognition and instruction following, enable typographic jailbreaks, and remain effective across multi-image, scene-coherent, or web-like presentations \citep{cheng2024typographic,gong2025figstep,wang2025multiimage,cao2025scenetap,qraitem2025webartifact}. Those studies establish attackability and often optimize attack strength. Our question is narrower and safety-directed: for an ordered damage label, how does a fixed source-level message change an initially correct decision when it is delivered through the image, text, or both? Modality-matched benign controls distinguish that change from the instability introduced by adding an overlay or prefix, and upward transitions provide a check against undirected label movement.

\paragraph{Prompt injection and cross-modal evidential conflict.}
Prompt injection asks whether instructions inside model-consumed content can override the trusted task instruction \citep{perez2022ignore,greshake2023indirect,zhan2024injecagent}; multimodal variants transmit such control through images or audio \citep{bagdasaryan2023abusing,downer2025text2vlm,nagaraja2025imageprompt}. A related literature examines evidential conflict: VLMs can favor misleading text over relevant images, with preferences varying by task and model \citep{deng2025wordsvision,pezeshkpour2025mixed,zhang2026misinformation}. Out-of-context misinformation work, by contrast, trains models to detect and explain false image--text pairings \citep{qi2024sniffer}. Our direct payloads behave as external instructions, whereas misleading payloads make a false low-damage claim without commanding a label. Comparing both against benign messages on the same sources connects these literatures to a specific failure mode: induced under-triage in disaster assessment.
